\documentclass[12pt]{article}
\usepackage[utf8]{inputenc}
\usepackage[margin=1in]{geometry}
\usepackage{amsmath}
\usepackage{graphicx}
\usepackage{booktabs}
\usepackage{hyperref}
\usepackage[round]{natbib}

\title{EEG-fMRI in Awake Rats Reveals Suppressed Whole-Brain Sensory Responsiveness During Absence Seizures, Corroborated by Mean-Field Simulations}
\author{Author A$^{1}$, Author B$^{1,2}$, Author C$^{1}$, Author D$^{2}$ \\
\small $^{1}$Grenoble Institut des Neurosciences, Universit\'{e} Grenoble Alpes, Grenoble, France \\
\small $^{2}$Inserm U1216, Grenoble, France}
\date{}

\begin{document}
\maketitle

\begin{abstract}
Absence epilepsy is characterized by transient lapses of consciousness during spike-and-wave discharges (SWDs), yet the neural mechanisms underlying impaired sensory processing during seizures remain poorly understood. Prior human neuroimaging studies have been conducted exclusively at rest, without controlled external stimulation. Here, we combined simultaneous EEG and zero-echo-time (ZTE) functional MRI in non-curarized awake GAERS rats ($n = 11$) with visual and whisker stimulation paradigms to directly measure hemodynamic responses to sensory input during ictal and interictal states. ZTE fMRI achieved 35.8 dB lower acoustic noise than standard echo-planar imaging (78.7 vs. 114.5 dB), enabling spontaneous seizure occurrence during scanning with minimal head motion (mean framewise displacement $0.69 \pm 0.32$ $\mu$m). We found that whole-brain hemodynamic responses to both visual and somatosensory stimulation were globally suppressed during ictal periods relative to interictal periods ($F$-contrast maps, $p < 0.05$ cluster-corrected). Cortical responses became negatively polarized during seizures despite ongoing stimulation, suggesting active suppression rather than mere absence of activation. A mean-field computational model based on the Adaptive Exponential (AdEx) integrate-and-fire neuron reproduced SWD dynamics and demonstrated restricted propagation of stimulation-evoked activity during the ictal regime, consistent with the fMRI observations. These results provide the first direct evidence that sensory processing is globally suppressed across the brain during absence seizures in awake, behaving animals.
\end{abstract}

\section{Introduction}

Absence epilepsy is a common form of generalized epilepsy characterized by brief episodes of impaired consciousness coinciding with bilateral, synchronous spike-and-wave discharges (SWDs) on electroencephalography (EEG) \citep{blumenfeld2005consciousness}. During these episodes, patients appear unresponsive to their environment, exhibiting behavioral arrest and disrupted perception, yet the neural basis of this sensory disconnection remains poorly characterized \citep{coenen2003absence}. Understanding how absence seizures disrupt sensory processing has important clinical implications, as these transient impairments contribute to attention deficits, learning difficulties, and other neuropsychiatric comorbidities associated with the condition.

Prior neuroimaging studies in humans have revealed complex patterns of cortical and thalamic activation during absence seizures, including increased thalamic BOLD signal and variable cortical responses \citep{bai2010default}. However, these studies have been conducted exclusively during rest, without controlled external stimulation, making it impossible to directly assess how the brain processes incoming sensory information during the ictal state. In rodent models, functional MRI studies of absence epilepsy have typically used curarized or anesthetized preparations \citep{nersesyan2004remote}, which fundamentally alter brain state and confound interpretation of seizure-related changes.

A major technical barrier to awake-state fMRI in rodents is the acoustic noise generated by rapid gradient switching in standard echo-planar imaging (EPI) sequences. Peak acoustic noise levels of approximately 115 dB produce significant stress in awake rats, increasing motion artifacts and suppressing spontaneous seizure generation---the very phenomenon under investigation \citep{depaulis2016gaers}. This challenge has limited the field to anesthetized or paralyzed preparations that preclude behavioral assessment of consciousness.

The Genetic Absence Epilepsy Rat from Strasbourg (GAERS) is a well-established model of absence epilepsy that exhibits spontaneous SWDs with behavioral arrest resembling human absence seizures \citep{depaulis2016gaers}. Electrophysiological studies have identified the somatosensory cortex (barrel field) as the likely cortical initiation zone for SWDs in this model, with subsequent engagement of thalamocortical loops \citep{meeren2002cortical, polack2007deep}. However, whether and how sensory-evoked responses are modulated during these seizures has not been investigated with whole-brain imaging.

Here, we address these gaps by establishing a zero-echo-time (ZTE) fMRI protocol that achieves dramatically reduced acoustic noise compared to EPI, enabling non-curarized awake imaging in GAERS rats with simultaneous EEG recording. We combined this technical advance with visual and whisker stimulation paradigms to directly compare hemodynamic responses to sensory input during interictal baseline and ictal states. Furthermore, we developed a mean-field computational model based on the Adaptive Exponential (AdEx) neuron to provide a mechanistic account of how SWD dynamics restrict the propagation of stimulus-evoked neural activity.

\section{Methods}

\subsection{Animals and Ethical Approval}

Eleven adult GAERS rats (8--12 months; 6 males, 5 females; $260 \pm 21$ g) were used. Animals were housed individually on a 12:12 h light-dark cycle at $22 \pm 2^\circ$C with 50--60\% humidity. All procedures were approved by the local ethics committee (Comit\'{e} Local GIN, C2EA-04) and complied with EU Directive 2010/63/EU.

\subsection{EEG Electrode Implantation}

Carbon fiber electrodes (WPI, Sarasota, FL) with approximately 5 mm brush-like tips and 30 mm leads were implanted under isoflurane anesthesia (5\% induction, 1--2\% maintenance). Recording electrodes were placed over right motor cortex (AP: $+2$ mm, ML: $+2.5$ mm) and right primary somatosensory cortex (AP: $-2.5$ mm, ML: $+3$ mm), with a cerebellar reference (AP: $-12$ mm, ML: $+2$ mm). Electrodes were secured with cyanoacrylate and dental cement (Selectaplus). Local anesthesia was provided with lidocaine hydrochloride.

\subsection{MRI Acquisition}

Imaging was performed on a 9.4T Bruker scanner with a custom transmit-receive loop coil (Table~\ref{tab:mri}). The ZTE sequence achieved dramatically lower acoustic noise than standard EPI while maintaining adequate sensitivity for functional imaging.

\begin{table}[htbp]
\centering
\caption{Comparison of MRI acquisition parameters for ZTE and EPI sequences.}
\label{tab:mri}
\begin{tabular}{lcc}
\toprule
Parameter & ZTE & EPI \\
\midrule
Peak acoustic noise (dB) & 78.7 & 114.5 \\
Noise difference & \multicolumn{2}{c}{35.8 dB ($\sim$62$\times$ weaker pressure)} \\
Spatial SNR & $\sim$25 & -- \\
Temporal SNR & 30--60 & -- \\
Signal change (visual stim.) & $\sim$0.3\% & 1--1.5\% \\
Susceptibility distortions & Minimal & Pronounced \\
EEG gradient artifacts & Low & High \\
\bottomrule
\end{tabular}
\end{table}

After a multi-week habituation protocol, animals underwent 45-minute awake EEG-fMRI sessions. Head motion was extremely low, with mean framewise displacement of $0.69 \pm 0.32$ $\mu$m and motion occurrences of $0.43 \pm 0.45$ per minute (Table~\ref{tab:motion}).

\begin{table}[htbp]
\centering
\caption{Motion parameters during awake ZTE fMRI.}
\label{tab:motion}
\begin{tabular}{lc}
\toprule
Metric & Value \\
\midrule
Mean framewise head translation ($\mu$m) & $0.69 \pm 0.32$ \\
Mean head rotation ($^\circ$) & $0.79 \pm 0.55$ \\
Maximum head displacement ($\mu$m) & $12.8 \pm 11.6$ \\
Maximum displacement (voxels) & $0.03 \pm 0.02$ \\
Motion occurrences/min (ZTE) & $0.43 \pm 0.45$ \\
Motion occurrences/min (EPI, prior) & $1.0 \pm 0.20$ \\
\bottomrule
\end{tabular}
\end{table}

\subsection{Sensory Stimulation}

Visual stimulation was delivered via bilateral optical fiber cables positioned near the eyes. Whisker stimulation was provided by bilateral air-puff tips directed at the whiskers. In each session, 6-second stimulation blocks were delivered during both interictal baseline periods and spontaneous absence seizures. EEG recordings were used to classify each stimulation block as ictal (delivered during SWD) or interictal post hoc.

\subsection{fMRI Analysis}

Whole-brain hemodynamic response functions (HRFs) were modeled using a general linear model (GLM) with gamma basis functions. Parameter estimates were computed voxel-wise, and $F$-contrast maps ($p < 0.05$, cluster-corrected) compared interictal versus ictal stimulation responses. ROI-based HRF analyses with extreme beta-value comparisons were performed for visual pathway structures (visual cortex, superior colliculus, lateral geniculate nucleus) and somatosensory pathway structures (barrel field cortex, ventral posteromedial thalamus). Activation maps were generated using Aedes software.

\subsection{Mean-Field Computational Model}

A mean-field simulation based on the Adaptive Exponential (AdEx) integrate-and-fire neuron model \citep{brette2005adaptive} was developed to reproduce both asynchronous irregular (AI) interictal dynamics and SWD ictal dynamics. The transition between states was controlled by modulating the adaptation current strength parameter. Virtual stimulation experiments tested whether activity propagation was restricted during the ictal regime compared to the interictal regime.

\section{Results}

\subsection{ZTE fMRI Enables Awake-State Imaging with Spontaneous Seizures}

The ZTE sequence produced peak acoustic noise of 78.7 dB, representing a 35.8 dB reduction compared to standard EPI (114.5 dB). This reduction corresponds to approximately 62-fold weaker sound pressure, enabling a stress-free awake state conducive to spontaneous seizure generation. Spontaneous SWDs at 2--5 Hz were reliably detected on simultaneous EEG during scanning sessions. The signal change with ZTE was approximately 0.3\%, lower than the 1--1.5\% typically observed with gradient-echo EPI, but ZTE provided substantially reduced susceptibility-induced distortions and EEG gradient artifacts. A pilot comparison in one animal confirmed that ZTE was superior to EPI for this application.

\subsection{Interictal Sensory Responses Follow Expected Pathway Activation}

During interictal visual stimulation, robust bilateral activation was observed in expected visual pathway structures including visual cortex, superior colliculus, and lateral geniculate nucleus of the thalamus ($F$-contrast maps, $p < 0.05$ cluster-corrected; Table~\ref{tab:visual}). Additionally, frontal cortex regions including prelimbic, cingulate, and secondary motor cortex showed significant positive activation, consistent with attentional processing of visual input.

\begin{table}[htbp]
\centering
\caption{Brain regions activated during interictal sensory stimulation.}
\label{tab:visual}
\begin{tabular}{llll}
\toprule
Stimulation & Brain Region & Activation & Pathway Role \\
\midrule
Visual & Visual cortex (bilateral) & Positive & Primary processing \\
Visual & Superior colliculus & Positive & Subcortical relay \\
Visual & Lateral geniculate nucleus & Positive & Thalamic relay \\
Visual & Frontal cortex (prelimbic) & Positive & Attention \\
Visual & Cingulate cortex & Positive & Executive processing \\
Whisker & Barrel field (S1) & Positive & Primary somatosensory \\
Whisker & VPM thalamus & Positive & Thalamic relay \\
Whisker & Frontal cortex & Positive & Attention \\
\bottomrule
\end{tabular}
\end{table}

During interictal whisker stimulation, strong positive activation was observed in the barrel field of somatosensory cortex and ventral posteromedial thalamus (VPM), with accompanying frontal cortex engagement. These activation patterns confirm that the awake ZTE fMRI protocol captures expected sensory-evoked hemodynamic responses with sufficient sensitivity.

\subsection{Global Suppression of Sensory Responses During Ictal Periods}

The central finding of this study is that hemodynamic responses to both visual and whisker stimulation were globally suppressed during ictal periods relative to interictal periods. $F$-contrast maps comparing interictal minus ictal responses revealed widespread significant differences across the respective sensory pathways ($p < 0.05$, cluster-corrected).

For visual stimulation, the visual cortex, superior colliculus, lateral geniculate nucleus, and frontal cortex all showed significantly reduced activation during seizures. Strikingly, cortical hemodynamic responses became negatively polarized during ictal visual stimulation, meaning that the BOLD signal deflected below baseline despite ongoing sensory input. This negative polarization suggests an active suppressive mechanism rather than a simple absence of activation.

The whisker stimulation group showed parallel findings: barrel field cortex, VPM thalamus, and frontal cortex all demonstrated suppressed responses during the ictal state. ROI-based HRF analyses confirmed that positive hemodynamic responses during interictal stimulation were replaced by negatively deflected responses during ictal stimulation, with extreme beta-value comparisons confirming statistically significant differences between states.

\subsection{Hemodynamic Signature of Seizures Without Stimulation}

Analysis of seizure-related hemodynamic changes in the absence of sensory stimulation revealed a complex pattern. The somatosensory cortex, believed to be the seizure initiation zone in GAERS, showed positive or mixed responses. Thalamic structures showed positive activation consistent with their essential role in sustaining SWDs. Other cortical regions showed variable responses reflecting seizure propagation patterns from the barrel cortex.

\subsection{Mean-Field Model Reproduces Restricted Activity Propagation}

The AdEx mean-field model successfully generated two distinct dynamical regimes by modulating adaptation current strength. At lower adaptation values, the network exhibited asynchronous irregular (AI) dynamics representing the interictal state. At higher adaptation values, the network transitioned to SWD dynamics, with the local field potential (LFP) and membrane potential traces capturing the characteristic alternating burst and silent periods observed experimentally.

Virtual stimulation experiments in the model revealed a critical difference between regimes. During the AI (interictal) state, stimulus-evoked activity propagated broadly to connected regions, consistent with normal sensory processing. During the SWD (ictal) state, stimulus-evoked activity remained confined to the stimulated region with markedly limited spatial spread, directly paralleling the fMRI observation of suppressed and spatially restricted sensory responses during seizures \citep{destexhe2009selforganized}.

\subsection{Summary of Findings Across Modalities}

The consistency of findings across visual and somatosensory modalities strengthens the conclusion that absence seizures produce a generalized suppression of sensory processing (Table~\ref{tab:summary}).

\begin{table}[htbp]
\centering
\caption{Summary of key findings across sensory modalities.}
\label{tab:summary}
\begin{tabular}{lcc}
\toprule
Finding & Visual & Whisker \\
\midrule
Interictal activation & Expected pathway & Expected pathway \\
Frontal cortex involvement & Present & Present \\
Ictal response suppression & Global, significant & Global, significant \\
Cortical HRF polarity (ictal) & Negative & Negative \\
Spatial extent (ictal vs. interictal) & Restricted & Restricted \\
Model consistency & Good & Good \\
\bottomrule
\end{tabular}
\end{table}

\section{Discussion}

This study provides the first direct demonstration that hemodynamic responses to controlled sensory stimulation are globally suppressed across the brain during absence seizures in awake, non-curarized animals. By combining a novel quiet ZTE fMRI protocol with simultaneous EEG in the GAERS model, we overcame the longstanding technical barrier of acoustic noise that has confined previous rodent epilepsy imaging to anesthetized or paralyzed preparations. The finding that cortical responses become negatively polarized during ictal stimulation, rather than simply being attenuated, suggests an active suppressive mechanism operating during SWDs.

\subsection{Technical Advance: Quiet fMRI for Awake Epilepsy Research}

The ZTE sequence used in this study achieved a 35.8 dB reduction in acoustic noise compared to standard EPI, representing an approximately 62-fold decrease in sound pressure. This reduction was critical for two reasons. First, the reduced noise minimized stress-related motion artifacts, as evidenced by mean framewise displacements below 1 $\mu$m. Second, and more importantly, the quiet environment preserved the calm awake state necessary for spontaneous seizure generation in GAERS rats, which are known to suppress SWDs under stress \citep{depaulis2016gaers}. Although ZTE produced smaller signal changes ($\sim$0.3\% versus 1--1.5\% for EPI), this was offset by the ability to collect data during natural seizure states. The ZTE contrast mechanism is sensitive to cerebral blood flow changes via T1-relaxation rate modulation, which has been suggested to be approximately 67\% more sensitive than standard BOLD EPI for detecting tissue oxygenation changes \citep{mishra2013zte, logothetis2001bold}.

\subsection{Active Suppression of Sensory Processing}

The negative polarization of cortical hemodynamic responses during ictal stimulation is a key finding that distinguishes active suppression from passive disengagement. If seizures merely diverted attentional resources away from sensory processing, one would expect attenuated but positive hemodynamic responses. Instead, the observation of negatively deflected responses suggests that the ongoing SWD dynamics actively suppress neural responses to incoming stimuli, potentially through enhanced inhibitory processes associated with the wave component of SWDs \citep{david2008imaging}.

This interpretation is supported by our mean-field model, which demonstrates that during the SWD regime, the network dynamics intrinsically restrict the propagation of externally evoked activity. The alternating excitatory bursts and inhibitory silent periods characteristic of SWDs create temporal windows during which incoming signals cannot effectively recruit downstream neurons, effectively confining stimulus-driven activation to the immediately stimulated region.

\subsection{Implications for Absence Epilepsy Pathophysiology}

Our results have direct relevance to understanding the consciousness impairment associated with absence seizures \citep{blumenfeld2005consciousness}. The global nature of sensory suppression---affecting both visual and somatosensory modalities, as well as frontal attentional networks---is consistent with the behavioral unresponsiveness observed during human absence seizures. The finding that the somatosensory cortex, which serves as the seizure initiation zone in GAERS, retains robust sensory responses during interictal periods confirms that the impairment is state-dependent rather than reflecting a permanent circuit abnormality.

The thalamic engagement observed during seizures without stimulation is consistent with the established role of thalamocortical loops in sustaining SWDs \citep{meeren2002cortical, polack2007deep}. The fact that thalamic sensory relay nuclei (LGN for vision, VPM for somatosensation) show suppressed responses during ictal stimulation suggests that the thalamic gating of sensory information is disrupted during seizures, potentially through aberrant feedforward inhibition and burst firing patterns \citep{bai2010default}.

\subsection{Limitations and Future Directions}

Several limitations should be considered. The ZTE signal change of approximately 0.3\% is substantially smaller than conventional EPI, potentially limiting sensitivity to weak or spatially restricted responses. The relatively small sample size ($n = 11$) is typical for awake animal fMRI studies but limits statistical power for detecting subtle between-subject differences. The mean-field model, while capturing the essential dynamics, is a simplification that does not account for the full complexity of thalamocortical circuitry. Future studies could employ cell-type-specific manipulations during awake imaging to dissect the contributions of specific neuronal populations to ictal sensory suppression.

\section{Conclusion}

By establishing quiet ZTE fMRI for awake epilepsy research, we have demonstrated that absence seizures produce global suppression of sensory hemodynamic responses across both visual and somatosensory pathways. The active nature of this suppression---evidenced by negatively polarized cortical responses---combined with computational modeling showing restricted activity propagation during SWD dynamics, provides a mechanistic framework for understanding the transient consciousness impairment that defines absence epilepsy. These findings bridge the gap between electrophysiological observations of SWDs and the functional disconnection from the sensory environment experienced by patients with absence seizures.

\bibliographystyle{plainnat}
\bibliography{references}

\end{document}
